%======================================================================
\section{Variables, Vectors, Matrices, and Lists}
%======================================================================

\sectionframe

\begin{frame}[fragile]{Operators \& Operator-like Functions}
\begin{itemize}
  \item \textbf{Variable Assignment} : The preferred method is using \code{<-} but \code{=} also works \textit{(Note: When you are in the R Coding environment, you can type \code{Alt + -} to enter \code{<-})}
\end{itemize}
\begin{lstlisting}[language=R]
x <- 5
y = 11
\end{lstlisting}
\begin{itemize}
  \item \textbf{Arithmetic Operators} : \code{+ - / *}
\end{itemize}
\begin{lstlisting}[language=R]
x + y 
x / y
x * y
\end{lstlisting}
\begin{lstlisting}[style=output]
[1] 16
[1] 0.4545455
[1] 55
\end{lstlisting}
\end{frame}

\begin{frame}[fragile]{Operators \& Operator-like Functions (cont.)}
\begin{itemize}
  \item \textbf{Sequences} : Use \code{:} to create a sequence of numbers.
  \item \textbf{Equality Comparison} : use \code{==} to test whether items are equal.
  \item \textbf{Inequality Comparison} : Use \code{< > <= >=} to test inequalities.
  \item \textbf{Negation} : Use \code{!} to negate a logical TRUE / FALSE value
\end{itemize}
\begin{lstlisting}[language=R]
1:10
5 == 6
5 <= 6
5 != 6
\end{lstlisting}
\begin{lstlisting}[style=output]
 [1]  1  2  3  4  5  6  7  8  9 10
[1] FALSE
[1] TRUE
[1] TRUE
\end{lstlisting}
\end{frame}

\begin{frame}[fragile]{Operators \& Operator-like Functions (cont.)}
\begin{itemize}
  \item \textbf{AND} : Use \code{\&}
  \item \textbf{OR} : Use \code{|}
  \item \textbf{To see if an element is contained in a vector/list} : Use \code{\%in\%}
  \item \textbf{TRUE/FALSE reserved words} : \code{TRUE, FALSE, T, F}
\end{itemize}
\begin{lstlisting}[language=R]
(5 == 5) & (4 < 3)
(5 == 5) | (4 < 3)
2 %in% 1:3
TRUE == T
FALSE == F
\end{lstlisting}
\begin{lstlisting}[style=output]
[1] FALSE
[1] TRUE
[1] TRUE
[1] TRUE
[1] TRUE
\end{lstlisting}
\end{frame}

\begin{frame}[fragile]{Operators \& Operator-like Functions}
\framesubtitle{Practice}
\begin{itemize}
  \item Do the following:
\end{itemize}
\begin{enumerate}
  \item Assign the following variables: 14 to a, 122 to b, 201 to c, and 5 to e.
\end{enumerate}
\begin{lstlisting}[language=R]
a <- 14
b <- 122
c <- 201
e <- 5
\end{lstlisting}
\begin{enumerate}\setcounter{enumi}{1}
  \item Assign the sum of a, b, and c to d
\end{enumerate}
\begin{lstlisting}[language=R]
d <- a + b + c
\end{lstlisting}
\begin{enumerate}\setcounter{enumi}{2}
  \item Assign the product of b times c to f
\end{enumerate}
\begin{lstlisting}[language=R]
f <- b * c
\end{lstlisting}
\end{frame}

\begin{frame}[fragile]{Operators \& Operator-like Functions (cont.)}
\framesubtitle{Practice}
\begin{enumerate}\setcounter{enumi}{3}
  \item Assign the sequence of numbers from 5 to 10 to g
\end{enumerate}
\begin{lstlisting}[language=R]
g <- 5:10
\end{lstlisting}
\begin{enumerate}\setcounter{enumi}{4}
  \item Determine which values in g are equal to 6
\end{enumerate}
\begin{lstlisting}[language=R]
g == 6
\end{lstlisting}
\begin{lstlisting}[style=output]
[1] FALSE  TRUE FALSE FALSE FALSE FALSE
\end{lstlisting}
\begin{enumerate}\setcounter{enumi}{5}
  \item Assign the sequence of numbers from 7 to 8 to h
\end{enumerate}
\begin{lstlisting}[language=R]
h <- 7:8
\end{lstlisting}
\end{frame}

\begin{frame}[fragile]{Operators \& Operator-like Functions (cont.)}
\framesubtitle{Practice}
\begin{enumerate}\setcounter{enumi}{6}
  \item Determine which values in g are greater than h
\end{enumerate}
\begin{lstlisting}[language=R]
g > h
\end{lstlisting}
\begin{lstlisting}[style=output]
[1] FALSE FALSE FALSE FALSE  TRUE  TRUE
\end{lstlisting}
\begin{enumerate}\setcounter{enumi}{7}
  \item Negate the results of Q5
\end{enumerate}
\begin{lstlisting}[language=R]
g != 6
\end{lstlisting}
\begin{lstlisting}[style=output]
[1]  TRUE FALSE  TRUE  TRUE  TRUE  TRUE
\end{lstlisting}
\begin{enumerate}\setcounter{enumi}{8}
  \item Negate the results of Q7
\end{enumerate}
\begin{lstlisting}[language=R]
!(g > h)
\end{lstlisting}
\begin{lstlisting}[style=output]
[1]  TRUE  TRUE  TRUE  TRUE FALSE FALSE
\end{lstlisting}
\end{frame}

\begin{frame}[fragile]{Operators \& Operator-like Functions (cont.)}
\framesubtitle{Practice}
\begin{enumerate}\setcounter{enumi}{9}
  \item Is a less than b and is b less than c?
\end{enumerate}
\begin{lstlisting}[language=R]
a < b & b < c
\end{lstlisting}
\begin{lstlisting}[style=output]
[1] TRUE
\end{lstlisting}
\begin{enumerate}\setcounter{enumi}{10}
  \item Is a greater than d or f?
\end{enumerate}
\begin{lstlisting}[language=R]
a > d | a > f
\end{lstlisting}
\begin{lstlisting}[style=output]
[1] FALSE
\end{lstlisting}
\begin{enumerate}\setcounter{enumi}{11}
  \item Is 5 in a,b,c,d,e, or f?
\end{enumerate}
\begin{lstlisting}[language=R]
5 %in% c(a,b,c,d,e,f)
\end{lstlisting}
\begin{lstlisting}[style=output]
[1] TRUE
\end{lstlisting}
\end{frame}

\begin{frame}{R Classes}
\begin{itemize}
  \item There are 4 R classes:
\end{itemize}
\begin{enumerate}
  \item \textbf{Vector} \textit{-- a one-dimensional data structure, where the elements are all of the same type}
  \item \textbf{Matrix} \textit{-- a two-dimensional table of vectors, where the elements are all of the same type}
  \item \textbf{List} \textit{-- a one-dimensional vector that can contain objects of different types}
  \item \textbf{Dataframe} \textit{-- a two-dimensional table of vectors, that can contain data of different types}
\end{enumerate}
\end{frame}

\begin{frame}[fragile]{Vectors}
\begin{itemize}
  \item \textit{A one-dimensional data structure, where the elements are all of the same type}
\end{itemize}
\textbf{There are 3 main ways to create a vector:}\par\smallskip
\begin{itemize}
  \item First, you could use \code{c()} function (c stands for “combine”):
\end{itemize}
\begin{lstlisting}[language=R]
v1 <- c(1, 2, 3, 4, 5)
v2 <- c("a", "b", "c")

# Not recommended
v3 <- c(1, "a")
\end{lstlisting}
\begin{columns}[T]
\begin{column}{0.48\textwidth}
\begin{lstlisting}[language=R]
v1
\end{lstlisting}
\begin{lstlisting}[style=output]
[1] 1 2 3 4 5
\end{lstlisting}
\end{column}
\begin{column}{0.48\textwidth}
\begin{lstlisting}[language=R]
v2
\end{lstlisting}
\begin{lstlisting}[style=output]
[1] "a" "b" "c"
\end{lstlisting}
\end{column}
\end{columns}
\end{frame}

\begin{frame}[fragile]{Vectors (cont.)}
\begin{lstlisting}[language=R]
v3
\end{lstlisting}
\begin{lstlisting}[style=output]
[1] "1" "a"
\end{lstlisting}
\begin{itemize}
  \item Second, you could use \code{seq(from, to, by)}:
\end{itemize}
\begin{lstlisting}[language=R]
v4 <- seq(1, 10, 2)
v5 <- seq(0, 1, 0.25)
\end{lstlisting}
\begin{columns}[T]
\begin{column}{0.48\textwidth}
\begin{lstlisting}[language=R]
v4
\end{lstlisting}
\begin{lstlisting}[style=output]
[1] 1 3 5 7 9
\end{lstlisting}
\end{column}
\begin{column}{0.48\textwidth}
\begin{lstlisting}[language=R]
v5
\end{lstlisting}
\begin{lstlisting}[style=output]
[1] 0.00 0.25 0.50 0.75 1.00
\end{lstlisting}
\end{column}
\end{columns}
\begin{itemize}
  \item Third, you could use \code{rep(elements, times)} to create a vector of repeated elements:
\end{itemize}
\begin{lstlisting}[language=R]
awesome <- rep("Boom", 4)
awesomer <- rep(c("three", "blind", "mice"), 2)
awesomest <- c(rep(NA,16),"Batman!")
\end{lstlisting}
\end{frame}

\begin{frame}[fragile]{Vectors (cont.)}
\begin{lstlisting}[language=R]
awesome
\end{lstlisting}
\begin{lstlisting}[style=output]
[1] "Boom" "Boom" "Boom" "Boom"
\end{lstlisting}
\begin{lstlisting}[language=R]
awesomer
\end{lstlisting}
\begin{lstlisting}[style=output]
[1] "three" "blind" "mice"  "three" "blind" "mice"
\end{lstlisting}
\begin{lstlisting}[language=R]
awesomest
\end{lstlisting}
\begin{lstlisting}[style=output]
 [1] NA        NA        NA        NA        NA        NA        NA       
 [8] NA        NA        NA        NA        NA        NA        NA       
[15] NA        NA        "Batman!"
\end{lstlisting}
\end{frame}

\begin{frame}[fragile]{Vectors}
\begin{itemize}
  \item \code{is.vector()} returns TRUE if an object is a vector:
\end{itemize}
\begin{lstlisting}[language=R]
is.vector(v5)
\end{lstlisting}
\begin{lstlisting}[style=output]
[1] TRUE
\end{lstlisting}
\begin{itemize}
  \item Use the \code{length()} function to obtain the length of a vector:
\end{itemize}
\begin{lstlisting}[language=R]
length(v1)
\end{lstlisting}
\begin{lstlisting}[style=output]
[1] 5
\end{lstlisting}
\begin{lstlisting}[language=R]
length(v2)
\end{lstlisting}
\begin{lstlisting}[style=output]
[1] 3
\end{lstlisting}
\end{frame}

\begin{frame}[fragile]{Vectors (cont.)}
\begin{itemize}
  \item We can do operations using vectors:
\end{itemize}
\begin{columns}[T]
\begin{column}{0.48\textwidth}
\begin{lstlisting}[language=R]
v1
\end{lstlisting}
\begin{lstlisting}[style=output]
[1] 1 2 3 4 5
\end{lstlisting}
\begin{lstlisting}[language=R]
v1 + 2
\end{lstlisting}
\begin{lstlisting}[style=output]
[1] 3 4 5 6 7
\end{lstlisting}
\end{column}
\begin{column}{0.48\textwidth}
\begin{lstlisting}[language=R]
# addition between two vectors
v1
\end{lstlisting}
\begin{lstlisting}[style=output]
[1] 1 2 3 4 5
\end{lstlisting}
\begin{lstlisting}[language=R]
v4
\end{lstlisting}
\begin{lstlisting}[style=output]
[1] 1 3 5 7 9
\end{lstlisting}
\begin{lstlisting}[language=R]
v1 + v4
\end{lstlisting}
\begin{lstlisting}[style=output]
[1]  2  5  8 11 14
\end{lstlisting}
\end{column}
\end{columns}
\end{frame}

\begin{frame}[fragile]{Vectors (cont.)}
\begin{lstlisting}[language=R]
# combining two vectors
v14 <- c(v1, v4)

# Adding two vectors that are not same length (Not recommended)
v6 <- c(1, 2, 3)

v7 <- c(1, 2, 3, 4, 5)

# Shorter: 1 2 3 1 2 (recycled)
# Longer:  1 2 3 4 5
# Result:  2 4 6 5 7
\end{lstlisting}
\begin{lstlisting}[language=R]
v6 + v7
\end{lstlisting}
\begin{lstlisting}[style=output]
Warning in v6 + v7: longer object length is not a multiple of shorter object
length
\end{lstlisting}
\begin{lstlisting}[style=output]
[1] 2 4 6 5 7
\end{lstlisting}
\end{frame}

\begin{frame}[fragile]{Vectors}
\begin{itemize}
  \item Extracting a subset of a vector:
\end{itemize}
\begin{lstlisting}[language=R]
vec1 <- c("e", "x", "a", "m", "p", "l", "e")
\end{lstlisting}
\begin{lstlisting}[language=R]
vec1[1]
\end{lstlisting}
\begin{lstlisting}[style=output]
[1] "e"
\end{lstlisting}
\begin{lstlisting}[language=R]
vec1[3:5]  # extract the 3rd to 5th elements
\end{lstlisting}
\begin{lstlisting}[style=output]
[1] "a" "m" "p"
\end{lstlisting}
\begin{lstlisting}[language=R]
vec1[c(4, 7)]  # extract the 4th and 7th elements
\end{lstlisting}
\begin{lstlisting}[style=output]
[1] "m" "e"
\end{lstlisting}
\begin{lstlisting}[language=R]
vec1[-1]  # negation indexing: extract all elements except the 1st
\end{lstlisting}
\begin{lstlisting}[style=output]
[1] "x" "a" "m" "p" "l" "e"
\end{lstlisting}
\end{frame}

\begin{frame}[fragile]{Vectors (cont.)}
\begin{lstlisting}[language=R]
vec1[-(3:5)]  # extract all elements except 3rd to the 5th elements
\end{lstlisting}
\begin{lstlisting}[style=output]
[1] "e" "p" "l" "e"
\end{lstlisting}
\begin{lstlisting}[language=R]
vec1[-c(4, 7)]  # extract all elements except the 4th and 7th
\end{lstlisting}
\begin{lstlisting}[style=output]
[1] "e" "x" "a" "p" "l"
\end{lstlisting}
\begin{itemize}
  \item Truncation using a logical vector:
\end{itemize}
\begin{lstlisting}[language=R]
x <- c(1, 10, 8, 2, 6)

x[x > 4]
\end{lstlisting}
\begin{lstlisting}[style=output]
[1] 10  8  6
\end{lstlisting}
\begin{lstlisting}[language=R]
# logical vector
x > 4
\end{lstlisting}
\begin{lstlisting}[style=output]
[1] FALSE  TRUE  TRUE FALSE  TRUE
\end{lstlisting}
\end{frame}

\begin{frame}[fragile]{Vectors (cont.)}
\begin{itemize}
  \item Override specific elements in a vector:
\end{itemize}
\begin{lstlisting}[language=R]
v4 <- rep(NA, 20)
v4
\end{lstlisting}
\begin{lstlisting}[style=output]
 [1] NA NA NA NA NA NA NA NA NA NA NA NA NA NA NA NA NA NA NA NA
\end{lstlisting}
\begin{lstlisting}[language=R]
v4[1:10] <- 1
v4
\end{lstlisting}
\begin{lstlisting}[style=output]
 [1]  1  1  1  1  1  1  1  1  1  1 NA NA NA NA NA NA NA NA NA NA
\end{lstlisting}
\begin{lstlisting}[language=R]
v4[11:20] <- 2
v4
\end{lstlisting}
\begin{lstlisting}[style=output]
 [1] 1 1 1 1 1 1 1 1 1 1 2 2 2 2 2 2 2 2 2 2
\end{lstlisting}
\end{frame}

\begin{frame}[fragile]{Vectors}
\framesubtitle{Practice}
\begin{itemize}
  \item Create the following vectors:
\end{itemize}
\begin{enumerate}
  \item Values from 55 to 75.
\end{enumerate}
\begin{lstlisting}[language=R]
seq(55,75)
\end{lstlisting}
\begin{lstlisting}[style=output]
 [1] 55 56 57 58 59 60 61 62 63 64 65 66 67 68 69 70 71 72 73 74 75
\end{lstlisting}
\begin{enumerate}\setcounter{enumi}{1}
  \item ‘You’, ‘are’, ‘the’, ‘best!’ repeated 3 times.
\end{enumerate}
\begin{lstlisting}[language=R]
rep(c('You', 'are',  'the', 'best!'),3)
\end{lstlisting}
\begin{lstlisting}[style=output]
 [1] "You"  "are"   "the"   "best!"  "You"  "are"   "the"   "best!"
 [9] "You"  "are"   "the"   "best!"
\end{lstlisting}
\end{frame}

\begin{frame}[fragile]{Vectors (cont.)}
\framesubtitle{Practice}
\begin{enumerate}\setcounter{enumi}{2}
  \item Even numbers from 20 to 84.
\end{enumerate}
\begin{lstlisting}[language=R]
seq(from = 20, to = 84, by = 2)
\end{lstlisting}
\begin{lstlisting}[style=output]
 [1] 20 22 24 26 28 30 32 34 36 38 40 42 44 46 48 50 52 54 56 58 60 62 64 66 68
[26] 70 72 74 76 78 80 82 84
\end{lstlisting}
\begin{enumerate}\setcounter{enumi}{3}
  \item What is the length of the vector created in Q3?
\end{enumerate}
\begin{lstlisting}[language=R]
length(seq(from = 20, to = 84, by = 2))
\end{lstlisting}
\begin{lstlisting}[style=output]
[1] 33
\end{lstlisting}
\begin{enumerate}\setcounter{enumi}{4}
  \item Assign 34 to y. Is y a vector?
\end{enumerate}
\begin{lstlisting}[language=R]
y <- 34
is.vector(y)
\end{lstlisting}
\begin{lstlisting}[style=output]
[1] TRUE
\end{lstlisting}
\end{frame}

\begin{frame}[fragile]{Matrix}
\begin{itemize}
  \item \textit{A two-dimensional table of vectors, where the elements are all of the same type}
  \item To create a matrix in R, we can use \code{matrix(data, [nrow], [ncol], [byrow])}:
\end{itemize}
\begin{lstlisting}[language=R]
A <- matrix(1:8)
\end{lstlisting}
\begin{columns}[T]
\begin{column}{0.3\textwidth}
\begin{lstlisting}[language=R]
A
\end{lstlisting}
\begin{lstlisting}[style=output]
     [,1]
[1,]    1
[2,]    2
[3,]    3
[4,]    4
[5,]    5
[6,]    6
[7,]    7
[8,]    8
\end{lstlisting}
\end{column}
\begin{column}{0.66\textwidth}
\begin{lstlisting}[language=R]
A_4x2 <- matrix(1:8, ncol = 2)
A_4x2
\end{lstlisting}
\begin{lstlisting}[style=output]
     [,1] [,2]
[1,]    1    5
[2,]    2    6
[3,]    3    7
[4,]    4    8
\end{lstlisting}
\end{column}
\end{columns}
\end{frame}

\begin{frame}[fragile]{Matrix (cont.)}
\begin{columns}[T]
\begin{column}{0.5\textwidth}
\begin{lstlisting}[language=R]
A_4x2_byrow <- matrix(1:8, ncol = 2,
                      byrow = TRUE)
A_4x2_byrow
\end{lstlisting}
\begin{lstlisting}[style=output]
     [,1] [,2]
[1,]    1    2
[2,]    3    4
[3,]    5    6
[4,]    7    8
\end{lstlisting}
\end{column}
\begin{column}{0.46\textwidth}
\begin{lstlisting}[language=R]
# the data can also be a vector
B <- matrix(c(2, 3, 1, 4),
            nrow = 2)
B
\end{lstlisting}
\begin{lstlisting}[style=output]
     [,1] [,2]
[1,]    2    1
[2,]    3    4
\end{lstlisting}
\end{column}
\end{columns}
\end{frame}

\begin{frame}[fragile]{Matrix (cont.)}
\begin{itemize}
  \item Looking at matrix attributes: \code{nrow()}, \code{ncol()}, \code{dim()}:
\end{itemize}
\begin{lstlisting}[language=R]
nrow(A_4x2)
\end{lstlisting}
\begin{lstlisting}[style=output]
[1] 4
\end{lstlisting}
\begin{columns}[T]
\begin{column}{0.48\textwidth}
\begin{lstlisting}[language=R]
ncol(A_4x2)
\end{lstlisting}
\begin{lstlisting}[style=output]
[1] 2
\end{lstlisting}
\end{column}
\begin{column}{0.48\textwidth}
\begin{lstlisting}[language=R]
dim(A_4x2)
\end{lstlisting}
\begin{lstlisting}[style=output]
[1] 4 2
\end{lstlisting}
\end{column}
\end{columns}
\begin{lstlisting}[language=R]
dim(A_4x2)[1]
\end{lstlisting}
\begin{lstlisting}[style=output]
[1] 4
\end{lstlisting}
\end{frame}

\begin{frame}[fragile]{Matrix}
\begin{itemize}
  \item Extract elements from a matrix using the indexing feature:
\end{itemize}
\begin{lstlisting}[language=R]
B
\end{lstlisting}
\begin{lstlisting}[style=output]
     [,1] [,2]
[1,]    2    1
[2,]    3    4
\end{lstlisting}
\begin{lstlisting}[language=R]
B[2, 2]  # second row & second column
\end{lstlisting}
\begin{lstlisting}[style=output]
[1] 4
\end{lstlisting}
\begin{columns}[T]
\begin{column}{0.48\textwidth}
\begin{lstlisting}[language=R]
B[1, ]  # first row & all columns
\end{lstlisting}
\begin{lstlisting}[style=output]
[1] 2 1
\end{lstlisting}
\end{column}
\begin{column}{0.48\textwidth}
\begin{lstlisting}[language=R]
B[, 2]  # all rows & second column
\end{lstlisting}
\begin{lstlisting}[style=output]
[1] 1 4
\end{lstlisting}
\end{column}
\end{columns}
\end{frame}

\begin{frame}[fragile]{Matrix (cont.)}
\begin{itemize}
  \item Basic operations using matrices:
\end{itemize}
\begin{lstlisting}[language=R]
C <- matrix(1:6, ncol = 2)
D <- matrix(7:12, ncol = 2)

C
\end{lstlisting}
\begin{lstlisting}[style=output]
     [,1] [,2]
[1,]    1    4
[2,]    2    5
[3,]    3    6
\end{lstlisting}
\begin{columns}[T]
\begin{column}{0.48\textwidth}
\begin{lstlisting}[language=R]
D
\end{lstlisting}
\begin{lstlisting}[style=output]
     [,1] [,2]
[1,]    7   10
[2,]    8   11
[3,]    9   12
\end{lstlisting}
\end{column}
\begin{column}{0.48\textwidth}
\begin{lstlisting}[language=R]
C + D
\end{lstlisting}
\begin{lstlisting}[style=output]
     [,1] [,2]
[1,]    8   14
[2,]   10   16
[3,]   12   18
\end{lstlisting}
\end{column}
\end{columns}
\end{frame}

\begin{frame}[fragile]{Matrix (cont.)}
\begin{lstlisting}[language=R]
C * D  # element-wise multiplication (NOT matrix multiplication)
\end{lstlisting}
\begin{lstlisting}[style=output]
     [,1] [,2]
[1,]    7   40
[2,]   16   55
[3,]   27   72
\end{lstlisting}
\begin{lstlisting}[language=R]
C == D  # to check equivalence, you must use "=="
\end{lstlisting}
\begin{lstlisting}[style=output]
      [,1]  [,2]
[1,] FALSE FALSE
[2,] FALSE FALSE
[3,] FALSE FALSE
\end{lstlisting}
\end{frame}

\begin{frame}[fragile]{Matrix (cont.)}
\begin{columns}[T]
\begin{column}{0.4\textwidth}
\begin{lstlisting}[language=R]
t(C)  # transpose
\end{lstlisting}
\begin{lstlisting}[style=output]
     [,1] [,2] [,3]
[1,]    1    2    3
[2,]    4    5    6
\end{lstlisting}
\end{column}
\begin{column}{0.56\textwidth}
\begin{lstlisting}[language=R]
C %*% t(D)  # matrix multiplication
\end{lstlisting}
\begin{lstlisting}[style=output]
     [,1] [,2] [,3]
[1,]   47   52   57
[2,]   64   71   78
[3,]   81   90   99
\end{lstlisting}
\end{column}
\end{columns}
\begin{lstlisting}[language=R]
solve(B)  # the inverse matrix
\end{lstlisting}
\begin{lstlisting}[style=output]
     [,1] [,2]
[1,]  0.8 -0.2
[2,] -0.6  0.4
\end{lstlisting}
\end{frame}

\begin{frame}[fragile]{Matrix}
\begin{itemize}
  \item Use \code{cbind()} to join by cols and \code{rbind()} to join by rows:
\end{itemize}
\begin{lstlisting}[language=R]
v1 <- c(1, 3, 5)
v2 <- c(2, 4, 6)

mat_r <- rbind(v1, v2) # bind these 2 vectors in a form of row vectors
\end{lstlisting}
\begin{lstlisting}[language=R]
mat_r
\end{lstlisting}
\begin{lstlisting}[style=output]
   [,1] [,2] [,3]
v1    1    3    5
v2    2    4    6
\end{lstlisting}
\begin{lstlisting}[language=R]
mat_c <- cbind(v1, v2) # bind these 2 vectors in a form of col vectors
\end{lstlisting}
\end{frame}

\begin{frame}[fragile]{Matrix (cont.)}
\begin{lstlisting}[language=R]
mat_c
\end{lstlisting}
\begin{lstlisting}[style=output]
     v1 v2
[1,]  1  2
[2,]  3  4
[3,]  5  6
\end{lstlisting}
\begin{itemize}
  \item The matrix will always be of one type. Therefore, if you are attempting to make a matrix with vectors of different types, R will coerce all the data to be of the most general type. For example:
\end{itemize}
\begin{lstlisting}[language=R]
a <- c(1,2,3)
b <- c(4,5,6)
c <- c("a","b","c")

my_matrix <- as.matrix(cbind(a,b,c)) # coerces all values to characters
\end{lstlisting}
\end{frame}

\begin{frame}[fragile]{Matrix (cont.)}
\begin{lstlisting}[language=R]
my_matrix
\end{lstlisting}
\begin{lstlisting}[style=output]
     a   b   c  
[1,] "1" "4" "a"
[2,] "2" "5" "b"
[3,] "3" "6" "c"
\end{lstlisting}
\begin{lstlisting}[language=R]
my_matrix2 <- as.matrix(cbind(a,b)) # keeps as numeric
\end{lstlisting}
\begin{lstlisting}[language=R]
my_matrix2
\end{lstlisting}
\begin{lstlisting}[style=output]
     a b
[1,] 1 4
[2,] 2 5
[3,] 3 6
\end{lstlisting}
\end{frame}

\begin{frame}[fragile]{List}
\begin{itemize}
  \item \textit{A one-dimensional vector that can contain objects of different types}
  \item Use \code{list()} to create a collection of components:
\end{itemize}
\begin{lstlisting}[language=R]
p <- "a list"
q <- 12:16
r <- matrix(1:9, ncol = 3)

# Putting objects p, q, and r into a list 
my_list <- list(p, q, r)
\end{lstlisting}
\end{frame}

\begin{frame}[fragile]{List (cont.)}
\begin{lstlisting}[language=R]
my_list
\end{lstlisting}
\begin{lstlisting}[style=output]
[[1]]
[1] "a list"

[[2]]
[1] 12 13 14 15 16

[[3]]
     [,1] [,2] [,3]
[1,]    1    4    7
[2,]    2    5    8
[3,]    3    6    9
\end{lstlisting}
\begin{itemize}
  \item Subset a list using square brackets:
\end{itemize}
\begin{lstlisting}[language=R]
my_list[[1]]
\end{lstlisting}
\begin{lstlisting}[style=output]
[1] "a list"
\end{lstlisting}
\end{frame}

\begin{frame}[fragile]{List (cont.)}
\begin{lstlisting}[language=R]
# extract a certain elements in a component
my_list[[2]][2] # extract the 2nd element in the 2nd component
\end{lstlisting}
\begin{lstlisting}[style=output]
[1] 13
\end{lstlisting}
\begin{itemize}
  \item Can also access subset by the names of component:
\end{itemize}
\begin{lstlisting}[language=R]
# Add names to components in a list (using "names()")
names(my_list) <- c("p", "q", "r")
my_list
\end{lstlisting}
\begin{lstlisting}[style=output]
$p
[1] "a list"

$q
[1] 12 13 14 15 16

$r
     [,1] [,2] [,3]
[1,]    1    4    7
[2,]    2    5    8
[3,]    3    6    9
\end{lstlisting}
\end{frame}

\begin{frame}[fragile]{List (cont.)}
\begin{columns}[T]
\begin{column}{0.48\textwidth}
\begin{lstlisting}[language=R]
# Access components by name in list
my_list$r
\end{lstlisting}
\begin{lstlisting}[style=output]
     [,1] [,2] [,3]
[1,]    1    4    7
[2,]    2    5    8
[3,]    3    6    9
\end{lstlisting}
\end{column}
\begin{column}{0.48\textwidth}
\begin{lstlisting}[language=R]
my_list$r[3, 3]
\end{lstlisting}
\begin{lstlisting}[style=output]
[1] 9
\end{lstlisting}
\end{column}
\end{columns}
\end{frame}

